\section{SI-Mapper : une approche agentique de la modélisation sémantique des bâtiments}
\label{sec:implementation}

La section précédente a établi que la cartographie manuelle des points BACnet est coûteuse, laborieuse
et résistante à l'automatisation avec les outils existants.
Cette section présente la solution proposée dans ce rapport, l'IA agentique pour la modélisation
sémantique HVAC automatisée, et retrace son développement complet depuis le concept initial jusqu'à
l'architecture finale, en passant par neuf itérations expérimentales.

\subsection{Vue d'ensemble de l'approche}
\label{subsec:si-mapper-overview}

SI-Mapper prend en entrée des données brutes de bâtiment, comme des plans d'ingénierie HVAC (PNG ou PDF) et
des exportations de points BACnet (fichiers CSV), et produit un modèle sémantique ASHRAE 223P au format
Turtle (TTL). L'ensemble du pipeline est orchestré par une application d'agent IA guidée par des instructions
de compétences spécifiques au domaine.
La figure~\ref{fig:pipeline} illustre le flux de données.

\begin{figure}[htbp]
\centering
\includegraphics[width=\textwidth]{pipeline.png}
\caption{Pipeline de données SI-Mapper --- des données brutes du bâtiment au modèle sémantique et à la visualisation en graphe.}
\label{fig:pipeline}
\end{figure}

Le facteur clé est l'IA multimodale.
Les grands modèles de langage récents comme Gemini 3.1 et Claude Sonnet 4.6 peuvent analyser directement
des plans d'ingénierie, en identifiant les symboles d'équipements, les connexions de gaines et les
relations spatiales à partir d'entrées visuelles.
Les schémas de systèmes HVAC, avec leur vocabulaire de symboles régulier et leurs conventions de mise en
page cohérentes, sont plus accessibles à la reconnaissance par IA que les plans d'étage architecturaux
généraux.
L'enseignement tiré du développement de SI-Mapper est que les modèles standard sont peu performants pour
la reconnaissance de plans spécifiques à un domaine jusqu'à ce qu'ils soient guidés par des instructions
détaillées et spécifiques au domaine, l'architecture de \emph{compétences} décrite dans la
section~\ref{subsec:skills}.

La nouveauté de SI-Mapper réside dans la combinaison de quatre capacités qu'aucun système précédent
n'intègre :
\begin{enumerate}[leftmargin=*, itemsep=2pt]
  \item \textbf{Reconnaissance multimodale de plans :} analyse par IA de plans d'ingénierie HVAC pour
        extraire l'équipement et la topologie des gaines.
  \item \textbf{Cartographie BACnet :} assignation automatisée des données de points BACnet aux équipements
        identifiés dans les plans.
  \item \textbf{Génération ASHRAE 223P :} production d'une ontologie ASHRAE 223P formellement valide
        incluant les ConnectionPoints, les Connections, les médiums et les liens de télémétrie.
  \item \textbf{Autocorrection itérative :} l'agent compare son modèle généré aux règles de validation,
        identifie les erreurs et les corrige de manière autonome --- en bouclant jusqu'à ce que le modèle
        soit valide.
\end{enumerate}

Cette approche remplace des semaines d'effort d'ingénierie expert par un pipeline automatisé qui s'exécute
en quelques minutes, produisant un modèle sémantique lisible par machine qui débloque la portabilité
logicielle, l'analytique avancée et l'intégration au réseau électrique pour le bâtiment.

% -----------------------------------------------------------------------
\subsection{Méthodologie de développement : neuf itérations expérimentales}
\label{subsec:iterations}
% -----------------------------------------------------------------------

Le développement de SI-Mapper a suivi une méthodologie expérimentale itérative.
Au fil de neuf itérations distinctes, l'architecture du système a évolué de simples scripts d'appel de
fonctions vers un système d'IA agentique sophistiqué basé sur les compétences. Chaque itération a été
motivée par les échecs et les limites de la précédente.

\begin{figure}[htbp]
  \centering
  \includegraphics[width=\textwidth]{experiment-progression.png}
  \caption{Progression des neuf itérations expérimentales --- des appels de fonctions aux compétences.}
  \label{fig:experiment-progression}
\end{figure}

\textbf{Itération 1 : LLMs comme fournisseurs de fonctions (Gemini 1.5 Flash).}
La première approche utilisait les LLMs comme fournisseurs d'intelligence intégrés dans des fonctions
Python conventionnelles : le modèle était appelé pour des tâches spécifiques telles que la reconnaissance
d'images ou la génération de texte, tandis que toute la logique métier restait codée en dur par le
développeur.
Cette approche était simple mais fondamentalement limitée, le LLM n'avait aucune conscience du contexte
applicatif et ne pouvait pas prendre de décisions autonomes. Chaque nouvelle capacité nécessitait
l'écriture de nouveau code autour d'elle. LangChain a été utilisé comme bibliothèque d'intégration à
ce stade.

\textbf{Itération 2 : Agents à base de flux de travail (Gemini 1.5 Flash).}
La deuxième itération a introduit des flux de travail pour définir le comportement de l'agent comme une
séquence d'étapes ordonnées, rendant possible une meilleure exploitation des capacités du LLM. Cependant,
la logique métier était toujours codée en dur dans les définitions de flux de travail, et l'agent restait
incapable de s'adapter à des situations inattendues. LangGraph a été évalué à ce stade mais a été
abandonné.

\textbf{Itération 3 : Flux de travail ADK (Gemini 2.0 Flash).}
La troisième itération a migré vers le Kit de Développement d'Agent (ADK) de Google tout en conservant
le modèle d'exécution basé sur les flux de travail. Un système RAG local a été introduit pour fournir aux
agents un accès à la demande à la documentation pour la classification des équipements HVAC. ADK venait
d'être publié à ce moment, avec une documentation et des exemples communautaires limités, ce qui signifiait
que les assistants de codage IA ne pouvaient pas aider de manière fiable avec la bibliothèque, une
limitation concrète qui a ralenti la vélocité de développement.

\textbf{Itération 4 : Agent maître agentique (Gemini 2.5 Flash).}
La quatrième itération représentait la première tentative agentique : un seul agent maître disposant d'une
pleine autonomie pour planifier et exécuter toutes les tâches. Cela a introduit une capacité de prise de
décision authentique mais a exposé un mode d'échec critique, les prompts nécessaires pour exprimer
le pipeline HVAC complet étaient trop longs et complexes pour que le modèle les suive de manière fiable.
Une simple boucle \texttt{while} ne suffisait pas à maintenir l'agent sur la bonne voie. Cette itération
a également introduit le frontend CopilotKit, le protocole Agent-to-UI pour l'observabilité en temps réel
de l'agent, et les outils MCP pour la gestion externe des outils.

\textbf{Itération 5 : Architecture Planifier-Agir-Réviser (Gemini 2.5 Flash).}
Inspirée par des modèles de conception multi-agents établis, la cinquième itération a décomposé
l'exécution en trois sous-agents séquentiels : Planificateur, Acteur et Réviseur. Cela a amélioré
la qualité de sortie mais l'agent terminait systématiquement les tâches avant d'avoir complété l'ensemble
du pipeline. L'ajout d'un agent de boucle sur la structure PAR a amélioré davantage les résultats mais
était encore insuffisant, l'agent ne pouvait pas compléter de manière fiable toutes les étapes en une
seule exécution.

\textbf{Itération 6 : Agent maître avec sous-agent PAR (Gemini 2.5 Flash).}
La sixième itération a ajouté un agent maître qui pouvait choisir d'agir directement ou de déléguer à
la structure PAR comme sous-agent. L'agent maître et le sous-agent disposaient tous deux d'une pleine
autonomie. Malgré cette flexibilité ajoutée, le sous-agent PAR ne pouvait toujours pas produire des
résultats d'une qualité suffisante. Les outils frontend ont été étendus avec le suivi des appels d'outils
et des variables d'état pour fournir une meilleure visibilité sur le raisonnement de l'agent.

\textbf{Itération 7 : Agent maître avec flux de travail hybrides (Gemini 2.5 Flash).}
La septième itération a remplacé l'agent PAR par une approche hybride : des flux de travail de
reconnaissance d'images spécialisés pour l'étape d'analyse des plans, combinés avec un planificateur
agentique pour l'orchestration des tâches. Cela a créé une architecture qui n'était ni entièrement basée
sur des flux de travail ni entièrement agentique. La qualité de la reconnaissance n'était pas suffisante,
et l'hybride architectural ajoutait de la complexité sans bénéfice proportionnel. Toute la séquence de
plans a finalement été reconvertie en flux de travail, mais les résultats restaient insuffisants.

\textbf{Itération 8 : Agents parallèles (Gemini 2.5 Flash, puis Gemini 3.0 Flash).}
La huitième itération a introduit l'exécution parallèle : un agent spécialisé par type d'équipement
s'exécutant simultanément, remplaçant les agents de reconnaissance séquentiels. Cela a significativement
amélioré la vitesse de traitement. Une mise à niveau du modèle vers Gemini 3.0 Flash a apporté une
amélioration de précision significative. Gemini 3.0 Pro a également été testé et a produit des résultats
de meilleure qualité mais nécessitait considérablement plus de temps.

\textbf{Itération 9 : Agent maître basé sur les compétences (Gemini 3.1 / Claude Sonnet 4.6).}
La neuvième et dernière itération a converti tous les sous-agents en compétences, des ensembles
d'instructions spécialisés exécutés en ligne dans le contexte de l'agent maître. Cette simplification
architecturale a été rendue possible par la capacité considérablement améliorée de suivi d'instructions
des derniers modèles. L'agent a également reçu la capacité de visiter le frontend via un navigateur sans
tête, de capturer une capture d'écran du canevas HVAC et de la comparer au plan source pour l'autocorrection
visuelle. L'architecture est devenue considérablement plus simple, plus facile à maintenir, et a produit
de meilleurs résultats que toute itération précédente. De plus, deux sous-agents exécutés en séquence ont
été maintenus uniquement pour créer le code Python et le code TTL du graphe. La raison de les maintenir
comme sous-agents plutôt que comme compétences est la complexité de la tâche qu'ils doivent exécuter.
Isoler le contexte pour l'exécution de la création du code dans des sous-agents était la meilleure
alternative.

L'enseignement clé de cette progression est que les avancées dans les capacités des modèles de fondation
ont rendu viables des architectures plus simples. L'approche basée sur les compétences n'est devenue
possible que lorsque des modèles comme Gemini 3.1 et Claude Sonnet 4.6 ont atteint une capacité suffisante
de suivi d'instructions pour exécuter des tâches de domaine complexes, à plusieurs étapes, dans un contexte
d'agent unique.

% -----------------------------------------------------------------------
\subsection{Architecture finale}
\label{subsec:final-architecture}
% -----------------------------------------------------------------------

\begin{figure}[htbp]
  \centering
  \includegraphics[width=\textwidth]{architecture.png}
  \caption{Architecture finale de SI-Mapper --- un agent maître unique avec des compétences spécialisées, des outils et une interface de visualisation.}
  \label{fig:architecture}
\end{figure}

L'architecture finale organise SI-Mapper en trois couches principales qui implémentent ensemble
le pipeline complet depuis les données multimodales brutes jusqu'à un modèle sémantique ASHRAE 223P
validé stocké dans une base de données de graphes.

\textbf{Frontend (CopilotKit + Next.js).}
Le frontend sert de tableau de bord opérationnel en temps réel. Un agent CopilotKit agit comme interface
de chat.
Le middleware dans ADK (callbacks avant et après le modèle) est utilisé pour diffuser les pensées, actions
et modifications d'état de l'agent au fur et à mesure dans le frontend, donnant aux opérateurs une
visibilité complète sur ce que fait l'agent et pourquoi.
Un canevas Graphivac restitue la grille HVAC interne en temps réel pendant que l'agent place gaines et
équipements à l'aide de ses outils.
Un visualiseur de code affiche le code Python ASHRAE 223P généré et le fichier TTL résultant.
Un visualiseur de graphe Sigma.js affiche le modèle sémantique final après son chargement dans Neo4j.

\textbf{Agent maître (Google ADK).}
Le cœur du système est un \textit{LlmAgent} maître encapsulé dans un agent de boucle, construit sur le
Kit de Développement d'Agent de Google, fonctionnant sur Gemini 3.1 ou Claude Sonnet 4.6. L'agent a accès
à un sous-agent séquentiel, environ 25 outils et six compétences.
Il maintient le contexte complet lors de l'extraction des données, du dessin du système HVAC et de la
cartographie des métadonnées BACnet.
L'agent maître délègue la création du code Python et la génération du TTL pour l'ontologie au sous-agent
séquentiel.

\textbf{Services externes.}
Deux services externes soutiennent le système : Neo4j et Graphivac. Neo4j avec le plugin Neosemantics (n10s)
stocke l'ontologie ASHRAE 223P comme un graphe RDF natif et fournit un accès aux requêtes Cypher
pour la traversée de graphe. L'API REST Graphivac gère le canevas CAD, permettant à l'agent
de lire et d'écrire l'état de la grille visuelle de manière bidirectionnelle.

\medskip
\noindent
Le choix d'une architecture à agent unique avec un seul sous-agent séquentiel a été motivé par plusieurs
avantages cumulatifs.
Les compétences fournissent une expertise de domaine en ligne, rendant le comportement de l'agent
entièrement observable dans une seule fenêtre de contexte sans l'opacité des appels de sous-agents
délégués.
L'agent maintient naturellement le contexte complet à travers toutes les étapes du pipeline, ce qui est
critique pour des opérations comme l'autocorrection visuelle, où l'agent doit comparer une capture
d'écran en direct du canevas avec le plan source et déterminer ce qui doit être corrigé.
Enfin, l'architecture est considérablement plus simple à maintenir, déboguer et étendre : ajouter une
nouvelle capacité signifie écrire une nouvelle compétence ou un nouvel outil, pas reconcevoir un protocole
de coordination d'agents.

% -----------------------------------------------------------------------
\subsection{Compétences}
\label{subsec:skills}
% -----------------------------------------------------------------------

Les compétences sont des ensembles d'instructions spécialisés que l'agent maître ou l'agent séquentiel
peuvent activer selon les besoins.
Ce ne sont pas des agents séparés, ils agissent comme un prompt dynamique.
Chaque compétence fournit des connaissances spécifiques au domaine pour une étape du pipeline, en
encodant les conventions de lecture de plans HVAC au niveau expert et les règles de génération d'ontologie
qu'un modèle de langage généraliste n'appliquerait pas de manière cohérente par lui-même.

\begin{table}[htbp]
  \caption{Compétences de l'agent maître et leurs rôles dans le pipeline.}
  \label{tab:skills}
  \centering
  \begin{tabular}{ll}
    \toprule
    \textbf{Compétence} & \textbf{Responsabilité} \\
    \midrule
    skill-ductwork            & Identifier et placer toutes les gaines des plans HVAC comme fondation spatiale \\
    skill-hvac-equipments     & Placer les équipements HVAC (ventilateurs, serpentins, volets, capteurs) sur le réseau de gaines \\
    skill-bacnet-points       & Extraire les données de points BACnet des exports CSV et mapper les points aux équipements \\
    skill-ontology-generation & Générer le code Python ASHRAE 223P à partir de l'état de la grille peuplée \\
    skill-ontology-validation & Corriger et valider itérativement le code d'ontologie pour produire un fichier TTL propre \\
    skill-read-code           & Charger les leçons de résolution d'erreurs des sessions de codage passées \\
    \bottomrule
  \end{tabular}
\end{table}

Les compétences suivent une chaîne de dépendances stricte qui reflète la réalité physique de la mise
en service HVAC. Les gaines sont placées en premier : tous les équipements doivent être ancrés à un
segment de gaine, donc la grille de gaines est la fondation spatiale sur laquelle tout le reste est
construit. Le placement des équipements suit, positionnant ventilateurs, serpentins, volets et capteurs
sur les coordonnées de gaines établies à l'étape précédente. La cartographie BACnet attache ensuite
les métadonnées de points en direct à chaque équipement de la grille. Chacune de ces trois étapes se
conclut par une étape de vérification humaine, permettant à l'opérateur de réviser et corriger le
travail de l'agent avant que le pipeline n'avance. La génération d'ontologie lit la grille entièrement
vérifiée et produit le code Python qui modélise le système en ASHRAE 223P. Enfin, la validation
d'ontologie entre dans une boucle d'autocorrection itérative, exécutant et corrigeant le code généré
jusqu'à ce qu'il produise un fichier TTL propre sans erreurs d'exécution. Chaque compétence suppose
que la précédente s'est terminée avec succès ; l'agent maître est responsable de la détection des
échecs et de la décision de réessayer ou d'escalader.

% -----------------------------------------------------------------------
\subsection{Outils}
\label{subsec:tools}
% -----------------------------------------------------------------------

L'agent maître a accès à environ 25 outils organisés en dix catégories.
Les outils gèrent toutes les interactions avec les systèmes externes, le stockage des données et l'état
opérationnel propre de l'agent, les compétences fournissent le raisonnement ; les outils fournissent
l'exécution.

\begin{table}[htbp]
  \caption{Outils de l'agent maître regroupés par catégorie.}
  \label{tab:tools}
  \centering
  \small
  \begin{tabular}{p{3.2cm}p{4.5cm}p{6.5cm}}
    \toprule
    \textbf{Catégorie} & \textbf{Outils} & \textbf{Fonction} \\
    \midrule
    Synchronisation grille
      & \texttt{sync\_graphivac\_to\_agent},
        \texttt{sync\_agent\_to\_graphivac}
      & Synchronisation bidirectionnelle entre l'état de l'agent et le canevas Graphivac \\
    \addlinespace
    CRUD grille
      & \texttt{add\_component},
        \texttt{add\_components\_batch},
        \texttt{delete\_component},
        \texttt{delete\_components\_batch},
        \texttt{read\_internal\_grid}
      & Créer, lire et supprimer des composants dans la grille interne \\
    \addlinespace
    Métadonnées
      & \texttt{write\_metadata},
        \texttt{write\_metadata\_batch}
      & Attacher les données de points BACnet aux équipements de la grille \\
    \addlinespace
    Ontologie
      & \texttt{search\_class\_mapping},
        \texttt{scan\_python\_files\_filtered},
        \texttt{read\_ontology},
        \texttt{write\_ontology},
        \texttt{execute\_ontology},
        \texttt{read\_prompt}
      & Recherche dans la bibliothèque ASHRAE 223P, lecture/écriture/exécution de code \\
    \addlinespace
    Signaux d'ontologie
      & \texttt{checkpoint\_code},
        \texttt{exit\_generator\_success},
        \texttt{exit\_generator\_failure},
        \texttt{exit\_validator\_success},
        \texttt{exit\_validator\_failure}
      & Versionnage des instantanés et signalisation de l'achèvement des tâches d'ontologie \\
    \addlinespace
    Neo4j
      & \texttt{load\_ttl\_to\_neo4j}
      & Importer le modèle sémantique TTL validé dans la base de données de graphes Neo4j \\
    \addlinespace
    Frontend
      & \texttt{capture\_frontend\_state}
      & Capturer une capture d'écran du canevas en direct pour l'autocorrection visuelle \\
    \addlinespace
    Ingestion
      & \texttt{ingest\_category\_files}
      & Charger les fichiers téléversés par l'utilisateur (plans, CSV, PDF) dans la session \\
    \addlinespace
    État / Progression
      & \texttt{save\_agent\_state},
        \texttt{get\_agent\_state},
        \texttt{update\_step},
        \texttt{update\_status}
      & Suivre la progression du traitement et persister l'état entre les tours \\
    \addlinespace
    Contrôle de boucle
      & \texttt{exit\_loop\_level\_2}
      & Terminer la session de l'agent maître \\
    \bottomrule
  \end{tabular}
\end{table}

% -----------------------------------------------------------------------
\subsection{Décisions technologiques clés}
\label{subsec:technology-decisions}
% -----------------------------------------------------------------------

Plusieurs choix technologiques ont été faits délibérément sur la base des leçons tirées du processus
de développement itératif. Chaque décision reflète un compromis entre capacité, coût et complexité
d'intégration.

\textbf{Google ADK.}
Le Kit de Développement d'Agent de Google a été choisi comme environnement d'exécution d'agent pour son
intégration native avec Gemini et son support intégré pour la gestion du cycle de vie des agents,
l'enregistrement des outils, l'état de session et la gestion des artefacts. Les itérations précédentes
utilisaient LangChain, mais les tests ont été réalisés aux premières étapes du développement de LangChain,
et à l'époque la bibliothèque avait encore de nombreux problèmes.
ADK a fourni depuis sa création l'infrastructure nécessaire pour réaliser les tests requis.
Cependant, au moment de la rédaction de ce rapport, de nouvelles alternatives apparaissent et d'autres
s'améliorent.
Ainsi, à ce stade, en considérant l'utilisation de LangChain, ce serait une alternative comparable à ADK.
De plus, l'introduction des agents profonds dans LangChain pourrait être une meilleure alternative.
De plus, tester de nouvelles capacités dans le cadre de Claude code est quelque chose qui doit être réalisé
à l'avenir pour avoir une vue d'ensemble de ce qui pourrait être le meilleur choix pour le cadre d'agent
pour continuer le développement.

\textbf{CopilotKit.}
CopilotKit implémente le protocole Agent-to-UI, permettant au backend Python ADK de diffuser
les pensées et actions de l'agent vers le frontend Next.js. Sans cette bibliothèque,
la construction de l'interface de chat aurait nécessité un développement supplémentaire sans valeur
de recherche. CopilotKit gère cette préoccupation, permettant au développement de se concentrer
sur le pipeline d'agent lui-même. Néanmoins, son utilisation pourrait être remise en question pour les
itérations futures.
Surtout si les itérations futures se concentrent sur le développement au-dessus de Claude code, auquel cas
une interface utilisateur pourrait être construite d'une manière complètement différente.

\textbf{Neo4j avec Neosemantics (n10s).}
ASHRAE 223P est une ontologie RDF ; son format de stockage naturel est une base de données de graphes avec
des capacités d'import RDF natives. Neo4j avec le plugin Neosemantics (n10s) satisfait les deux
exigences : il stocke le fichier TTL comme un graphe natif, et il fournit Cypher, un
langage de requête mature et expressif, pour la traversée et l'analyse de graphes. Des approches
alternatives comme le stockage TTL sous forme de fichiers plats ou l'utilisation d'un triple store ont
été envisagées, mais la combinaison de Neo4j de stockage natif de graphe et d'outillage ecosystémique
solide en a fait le choix pragmatique. Neo4j présente l'inconvénient de nécessiter la conversion du TTL
en un graphe ATTIC, et certaines informations pourraient être perdues dans cette conversion.
Les informations perdues n'ont pas été analysées dans cette itération en raison de limitations de temps.
Cependant, il est important de les analyser à l'avenir pour mieux comprendre les implications de la
perte d'informations.
GraphDB n'a pas le même problème car il ne nécessite pas de conversion depuis le TTL.
Cependant, GraphDB fonctionne avec SPARQL au lieu de Cypher.
Des travaux récents ont démontré que les LLMs atteignent une précision plus élevée lors de l'écriture de
requêtes Cypher que lors de l'écriture de requêtes SPARQL~\cite{Sivasubramaniam2024}.

\textbf{Sigma.js pour la visualisation de graphes.}
Sigma.js restitue de grands graphes en utilisant WebGL, offrant des performances fluides sur des modèles
sémantiques avec des centaines de nœuds et de relations. Des alternatives plus lourdes comme D3 ou Gephi
ont été envisagées, mais les deux introduisent une complexité inutile pour une exigence de visualisation
native dans le navigateur. Sigma.js s'intègre proprement avec le frontend Next.js et ne nécessite aucune
accommodation de rendu côté serveur au-delà d'un garde d'import dynamique standard.
